% 第 34 章　镜子、透镜与眼睛
\chapter{镜子、透镜与眼睛}

\begin{mainline}

\step{反常识开场}

先做一个你做过一万次、却未必解释得清的动作：站到镜子前，举起你的左手。

镜子里那个人举起了手。问题是：他举的是哪只手？几乎所有人的第一反应都是“右手”——镜子把左右颠倒了。可是再想想：镜子为什么不把上下也颠倒？镜面是平的，水平方向和竖直方向对它来说没有任何区别，它凭什么偏偏挑“左右”下手？如果镜子真的偏爱某个方向，那你把镜子横过来躺在地上照，是不是就该变成上下颠倒了？你可以回家试，结果不会让你满意：不管镜子怎么放，被“颠倒”的永远是左右，从来不是上下。

再看第二个现象。很多理发店的镜子只有半人高，挂得也不高，可你后退一两步，照样能在里面看到自己的全身，从头顶到鞋底。直觉说：镜子只有一半高，最多只能装下一半的你吧？事实恰好相反，而且——这是本章会让你反复吃亏的地方——你再往后退十步，需要的镜子高度一毫米都不会增加。

还有一个更日常的悬念。超市货架转角处挂着一面鼓鼓的凸面镜，汽车右侧后视镜上印着一行小字：“镜中物体比看起来更近”。为什么工程师故意用一面“说谎”的镜子？它说的一定是谎话吗？

这三个现象用的都是同一块物理：光走直线，遇到界面改变方向，而我们的眼睛和大脑只会按“光是直线来的”去倒推物体的位置。本章的任务，就是把这块物理拆开，顺便说清楚你的眼睛、眼镜、相机、显微镜和望远镜，以及——一件更要紧的事——这套“光线”理论在哪里会整个垮掉。

顺带一提，同样的倒推也在水面发生。湖边无风时，对岸的山和树在水里有一幅分毫不差的倒影——水面就是一面平放的平面镜，虚像悬在水面之下，与实景等深。风一起，波纹把镜面拆成无数朝不同方向倾斜的小镜片，倒影就碎成了晃动的光带。倒影的“清晰程度”，从头到尾只是一份关于水面平整度的报告。

\step{先猜一猜}

在往下读之前，先把你的直觉押在桌面上，不许待会儿偷偷改：

\begin{itemize}
  \item 猜一：镜子里的像“左右颠倒”，是因为镜面在水平方向做了什么竖直方向没做的事。你同意吗？如果同意，它做的是什么？
  \item 猜二：一个身高 1.7 米的人，想在一面竖直平面镜里看到全身，镜子至少要多高？和你的身高有什么关系？和你站多远有没有关系？
  \item 猜三：拿一片凸透镜（放大镜）对着窗外的楼，在透镜后面的白纸上，你猜会看到什么？正立的还是倒立的？如果你用手遮住透镜的左半边，纸上的像会发生什么——左半边消失？整个变暗？还是别的？
  \item 猜四：近视眼不戴眼镜时看不清远处。你猜他的眼球是“太会聚光”了，还是“不太会聚光”？
\end{itemize}

把答案写下来。本章结束时你要回来批改自己的卷子。

\step{动手实验}

\begin{experiment}[半面镜子、一把汤匙和一张纸]

\textbf{实验一：半身镜。} 找一面全身镜，背对窗户站好，让同伴拿两张便利贴：一张贴在镜面上正好挡住你头顶反射光的位置（头顶在镜中像的上方边缘），一张贴在正好挡住脚底的位置。量一量两张便利贴之间的距离，再量你的身高。然后你向前走近两步、再退后三步，请同伴重新贴一次。你会发现：便利贴的位置几乎不动，间距始终是身高的一半左右。哪怕退到房间另一头，结论也不变。这是平面镜成像规律递给你的第一份情报。

\textbf{实验二：一把汤匙。} 拿一把不锈钢汤匙。先看凸出的背面：你的脸是正立的、缩小的，怎么调距离都是正立缩小。再翻过来，看凹进去的一面：凑得很近时，脸是正立放大的；慢慢把汤匙拿远，某个瞬间像会糊成一片，再远，脸忽然翻成倒立的，而且越拿远缩得越小。同一块曲面，两种行为，中间隔着一个“翻车的距离”。这个距离就是本章的核心新词：焦距。

\textbf{实验三：墙上的风景。} 白天拉上窗帘，只留一条缝，让窗外明亮的景物（楼、树、路灯）能照亮房间一角。拿一片凸透镜——老花镜、放大镜都行——对着窗外，在透镜后方几十厘米处竖一张白纸，前后移动白纸。某个位置上，纸上会出现一幅窗外的风景：清晰、彩色，而且上下左右全部颠倒。这幅“画”不依附于任何颜料，你把纸拿走，它还在空气里；你伸手到那个位置，光斑落在手上。光线是真的在那里会聚了。这就是实像，是本实验要亲手抓住的东西。顺手做猜三的另一半：用手遮住透镜的一半。像没有缺一半，它只是变暗了。记下这个反直觉的结果，后面要解释。

\end{experiment}

\step{旧模型遇到麻烦}

关于“看见”，人类有过一个寿命极长的旧模型：眼睛里伸出某种“视线”，像手一样摸到物体，于是你看见了它。古希腊不少学者真这么想。这个模型有一个朴素的好处——解释了“我看哪儿就见哪儿”的主动性。但它的麻烦也朴素：漆黑的房间里，你的眼睛没有任何毛病，却什么也摸不到；捂住眼睛再松开，看见的过程不消耗任何可觉察的时间，这只“手”伸得未免太快。更要命的是镜子：镜中的椅子明明不在镜面后面，你的“视线手”却摸到了它——它什么时候学会穿墙了？

把旧模型换成第 33 章建立的：物体向四面八方发出或反射光，其中一小束钻进瞳孔，你才看见它。这个模型立刻解释了黑暗，但它自己也有个半成品问题——\emph{像是住在哪里的一幅画吗？}很多人下意识把镜中像当成“贴在镜面上的一张图”。这个模型同样一触即溃：你向左走，镜中像的位置跟着变；而站在另一侧的朋友同时看这面镜子，他看到的像和你看到的在空间的同一个深处——两个人从两个方向看“同一张贴在镜面的画”，不可能得到这个一致的纵深。像不在镜面上。像甚至不“在”任何地方等着被看；它是光线到达你眼睛之前排好的一个局。

还有一个更深一层的麻烦，先埋在这里：如果光严格走直线，那么墙上那幅实像的每个像点，应该只由一条光线决定。可实验三里，遮住一半透镜，像只是变暗——说明每个像点是由\emph{许多条}光线共同堆出来的。“一条光线对应一个点”的图像，马上也要升级。

\step{发明新概念}

我们现在像十九世纪的仪器师一样，把需要的概念一件件造出来。

\textbf{第一件：物点与光束。} 把物体拆成许多点，每个点向各个方向发光。照进你眼睛、照到镜面、穿过透镜的，永远只是这个点发出的一个细细的光锥。要追踪“看见”，只需追踪这个光锥被器件改造成了什么样。

\textbf{第二件：实像与虚像。} 一个物点发出的光锥，经过镜子或透镜之后，如果重新会聚到一个点，光真的在那里相交，那里就是\emph{实像点}——实验三里落在纸上的风景就是无数实像点的集合。如果光锥被弄得更发散，但这些光线反着往回延长时恰好交于一点，那里就是\emph{虚像点}：光没去过那里，可你的眼睛无法分辨“光是从那里发散出来的”和“光是被折腾成好像从那里发散出来的”。眼睛和大脑只认进入瞳孔的光锥形状，于是你“看见”了一个不在那里的东西。平面镜里的你，就是虚像。这不是幻觉的“假”，而是光路意义上的“不在场”。

\textbf{第三件：焦点与焦距。} 一束平行于主轴的光射向凹面镜或凸透镜，会聚后经过（或看起来经过）轴上的同一个点，叫\emph{焦点}；焦点到器件的距离叫\emph{焦距}。焦距是一个器件的身份证：它告诉你这东西把光“掰弯”的本事有多大。对球面镜，在近轴光线下焦距等于球面半径的一半；对透镜，焦距由两个球面的弯曲程度和玻璃折射率决定。汤匙实验里“像翻车”的那个距离，就是凹面的焦距：物点恰好落在焦点上时，反射光变成平行光，像被推到无穷远，于是糊成一片。

\textbf{第四件：三条特征光线。} 要预言像在哪里，不必追光锥里所有光线，只要三条好算的：平行于主轴入射的，过焦点；过焦点入射的，出来平行于主轴；直射透镜中心（或球面镜顶点）的，方向几乎不变。两条定交点，第三条用来检查。下面两幅图是本章的“主力图纸”，请对照着读。

有了焦点和焦距，球面镜家族在日常里的分工就一目了然了。凹面镜会聚：太阳灶和卫星接收天线把大面积收集来的平行信号聚到焦点的一个小锅上；手电筒和车灯反着用，把焦点处的光源变成平行光束射出去；化妆镜利用 $s<f$ 的正立放大虚像，把脸放得很大——代价是你必须凑到焦距以内。凸面镜发散：视野大、像缩小，于是它去看守超市转角、车库出口和汽车的右侧。还有一对容易混的邻居：房门上的“猫眼”是一片凹透镜，让屋内人看到门外正立缩小的大视野虚像；牙医用的小镜子则只是平面镜，负责拐个弯，不负责放大。

\begin{center}
\begin{tikzpicture}[scale=1.0,>=Stealth]
  \draw[thick] (0,-2.1)--(0,2.1);
  \fill[pattern=north east lines] (0,-2.1) rectangle (0.18,2.1);
  \node[right] at (0.1,-2.35) {镜面};
  \draw[->,very thick] (-3,0)--(-3,1.5);
  \node[left] at (-3,0.75) {物};
  \fill (-3,1.5) circle (1.2pt) node[above left]{$P$};
  \draw[->,very thick,dashed] (3,0)--(3,1.5);
  \node[right] at (3,0.75) {虚像};
  \fill (3,1.5) circle (1.2pt) node[above right]{$P'$};
  \draw[->] (-3,1.5)--(0,0.9);
  \draw[->] (0,0.9)--(-1.5,-0.6);
  \draw[->] (-3,1.5)--(0,0.1);
  \draw[->] (0,0.1)--(-1.5,-0.6);
  \draw[dashed] (0,0.9)--(3,1.5);
  \draw[dashed] (0,0.1)--(3,1.5);
  \draw (-1.5,-0.6) circle (0.14);
  \node[below] at (-1.5,-0.85) {眼睛};
  \draw[dashed,gray] (-3,0)--(3,0);
  \node[below] at (0,-3.0) {图 1：平面镜。反射光线反向延长交于镜后的 $P'$，像与物关于镜面对称。};
\end{tikzpicture}
\end{center}

\begin{center}
\begin{tikzpicture}[scale=0.78,>=Stealth]
  \draw[->] (-7.2,0)--(7.2,0);
  \draw[thick] (0,-2.3)--(0,2.6);
  \draw[->,thick] (0,2.6)--(0.18,2.42);
  \draw[->,thick] (0,-2.3)--(-0.18,-2.12);
  \draw[->,thick] (0,2.6)--(-0.18,2.42);
  \draw[->,thick] (0,-2.3)--(0.18,-2.12);
  \fill (2,0) circle (1pt) node[below]{$F$};
  \fill (4,0) circle (1pt) node[below]{$2F$};
  \fill (-2,0) circle (1pt) node[below]{$F$};
  \fill (-4,0) circle (1pt) node[below]{$2F$};
  \draw[->,very thick] (-5,0)--(-5,1.6);
  \node[left] at (-5,0.8) {物};
  \draw[->] (-5,1.6)--(0,1.6);
  \draw[->] (0,1.6)--(2,0);
  \draw (2,0)--(3.33,-1.07);
  \draw[->] (-5,1.6)--(0,0.53);
  \draw (0,0.53)--(3.33,-1.07);
  \draw[->,very thick] (3.33,0)--(3.33,-1.07);
  \node[right] at (3.33,-0.55) {实像};
  \fill (3.33,-1.07) circle (1.2pt);
  \node[below] at (0,-2.9) {图 2：凸透镜。物在二倍焦距外，像在 $F$ 与 $2F$ 之间，倒立缩小。};
\end{tikzpicture}
\end{center}

\textbf{第五件：眼球是一台什么机器。} 现在拆你自己的眼睛。角膜加晶状体相当于一组透镜，把外界成像在视网膜上；视网膜是感光面。这里用得上“眼球像相机”的比喻，但请先听完声明。\emph{它像什么}：前面有可会聚光的透镜组，后面有感光面，透镜到感光面的距离固定，靠改变晶状体的弯曲程度（改变焦距）来对焦，这一点和镜头伸缩对焦异曲同工；成的像一样是倒立缩小的实像。\emph{它不像什么}：视网膜上没有均匀的“像素网格”，中心一小片区域（黄斑）挤满了感光细胞，边缘稀疏得只能察觉动静；眼睛从不静止地“拍一张”，而是不断跳动扫描，你自以为看到的稳定高清世界，大半是大脑拿零碎素材补全的。把眼睛当相机用，学到光学；把眼睛当相机信，会学错神经。

这里立刻蹦出一个著名的问题：视网膜上的像是倒立的，为什么我们看到的世界是正的？看清本章的概念后，这个问题自己就散了：视网膜上根本没有一个“需要被谁再看一眼的小人”。感光细胞只负责把落在自己身上的光强和颜色转成电信号送出去，大脑拿到的是一堆编号好的信号线，而不是一张需要“摆正”的照片。“上”和“下”是大脑根据重力、身体姿态和手的反馈，给这些信号线贴的标签，贴标签的规则完全由后天校准。心理学家让人连续戴多天左右反转、上下颠倒的眼镜，起初世界天翻地覆，几天后大脑重新贴好标签，骑车倒水如常；摘下眼镜，又要重新适应一次。光学把像成在哪里、倒还是正，是物理；你觉得世界“正”，是标注。两者各管一段，互不拖欠。

\end{mainline}

\begin{mainline}

\step{一句话模型}

\begin{oneline}
镜子和透镜从不“生产”图像，它们只按反射定律和折射定律把每条光线换一个方向；同一个物点发出的光被改道之后若重新交于一点，那里就是实像，若只是反向延长线交于一点，那里就是虚像——你的眼睛只认光锥，不认产地。
\end{oneline}

\step{数学翻译器}

直觉已经就位，现在给它配一台计算器。全部成像问题压缩成一个方程加一条放大率。先立规矩（符号约定）：物在入射光一侧，物距 $s$ 为正；实像（光真的到达的那一侧）像距 $s'$ 为正，虚像为负；会聚器件（凸透镜、凹面镜）焦距 $f$ 为正，发散器件（凹透镜、凸面镜）焦距为负。在这套约定下，薄透镜和球面镜共用同一个式子：

\[
\frac{1}{s}+\frac{1}{s'}=\frac{1}{f},
\qquad
m=-\frac{s'}{s}.
\]

左边要回答的问题是“物放在 $s$，像出现在哪儿、多大”；右边唯一的输入 $f$ 是器件的会聚本事。$m$ 是像高与物高之比，负号代表倒立。

按惯例，公式落地先做极端体检：

\begin{itemize}
  \item \textbf{$s\to\infty$}：$1/s\to0$，得 $s'=f$。平行光聚在焦点——这正是焦距的定义，公式把定义吐了回来，没有闹独立。用放大镜聚焦阳光烤出焦斑，就是利用这一条：太阳在无穷远，光斑落在焦平面上。
  \item \textbf{$s=2f$}：$1/s'=1/f-1/(2f)=1/(2f)$，$s'=2f$，$m=-1$。等大倒立，物像对称。图 2 里物再往远挪，像就往 $F$ 缩，这正是相机拍远景时镜头几乎停在焦平面上的原因。
  \item \textbf{$s\to f$}：$1/s'\to0$，$s'\to\infty$。物在焦点，像被推到无穷远——出射光变平行。把这个方向反过来用，把灯泡放在焦点上，就得到探照灯和手电筒的抛物面反光碗。
  \item \textbf{$s<f$}：$1/s'=1/f-1/s<0$，$s'$ 为负——虚像，且 $|m|>1$，正立放大。这就是放大镜：报纸上每个字发出的光穿过透镜后变得更发散，你的眼睛倒推出一个更远更大的虚字。
  \item \textbf{$f<0$（凹透镜、凸面镜）}：对任意正的 $s$，$s'$ 恒为负，$0<m<1$——永远正立缩小虚像。凸面镜和汽车右后视镜把大片视野压缩进小镜面，代价是像变小、显得更远，所以镜子上才印着那句警告。它不是“说谎”，它是用缩小换视野，你需要重新校准距离感。
  \item \textbf{$f\to\infty$}：$1/f\to0$，$s'=-s$，$m=+1$。一个完全不会聚光的“透镜”就是平面镜：像与物等距、等大、正立、虚像。平面镜被干净地收编为同一公式的一个极限，这是好理论的标志——旧结论没有另立门户。
\end{itemize}

现在做一笔带真实数字的账，标的物是你此刻正在用的眼睛。人眼从晶状体到视网膜的距离约 $1.7$ 厘米，这是固定的 $s'$。看很远处（$s\to\infty$）时，晶状体放松，$f=1.7$ 厘米。看 $25$ 厘米处的书时，$1/f=1/0.25+1/0.017$，算出 $f\approx1.6$ 厘米——睫状肌把晶状体捏得更凸，焦距缩短了约 $1$ 毫米，就完成了从远山到书本的对焦。这个调节量小得惊人，也正是它撑不住时会出事：近视眼的眼球前后径偏长，远处物体成像在视网膜\emph{前方}（猜四的答案：不是不会聚，是太会聚了），于是配凹透镜先把光发散一点。眼镜度数用屈光度，$P=1/f$（$f$ 以米计）：一副 $-2.00$ 屈光度的镜片焦距 $-0.5$ 米，它把无穷远的物体成一个虚像在眼前 $0.5$ 米处——恰好落在近视者能看清的范围内。远视和老花则相反，是近处物体的像落到视网膜\emph{后方}，用凸透镜补一把会聚。

\textbf{把器件组合起来：相机、显微镜、望远镜。} 单个透镜的公式，已经够把三大光学仪器看穿。

相机是三者里最诚实的：它就是图 2 的直接应用。拍远处时 $s'\approx f$，感光元件就贴在焦平面附近；拍近处时 $s$ 变小，$s'$ 必须拉长，镜头于是“伸出去”一截，这就是对焦马达在做的事。镜头里还有个可变孔径叫光圈，用 $f$ 数（焦距除以孔径）标记：开大光圈，进光多，光锥张角大，焦外的点被糊成大斑，于是背景虚化、主体突出；收小光圈，前后景物都落在可容忍的清晰范围内——但收得太过（比如 $f/16$ 以下），每个点光源的衍射斑开始比像素还大，整个画面反而变软。选光圈就是“进光量、清晰纵深、衍射模糊”三方的谈判，这道谈判题用到本章全部三道围墙。

显微镜是两级接力。物镜焦距极短（毫米量级），把贴在焦距外侧一点点的标本成一个放大的倒立实像；目镜再把这个实像当“物”，按放大镜的老办法（$s<f$，正立放大虚像）送进眼睛，总放大率是两级相乘。这解释了显微镜物镜为什么又短又粗：短焦距意味着标本几乎贴着焦点，像距却拉到十几厘米，$m=-s'/s$ 自然大。几百倍的放大听起来唬人，其实每一步都只是那个一元方程在值班。

望远镜对付的是另一类问题：天体在无穷远，放大的不是“大小”而是“张角”。物镜焦距长，把遥远天体在焦平面上成一个实像；目镜焦距短，把这个实像当放大镜的物来看。角放大率 $M=f_{\text{物}}/f_{\text{目}}$，换一支短焦距目镜就换一种放大率，所以玩天文的人手里总攥着一排目镜。但请记住分辨率那笔账：放大只是把同一个斑放得更大，决定能不能看清细节的是口径 $D$。放大倍数超出口径供养能力的部分叫“空放大”——像越放越大，细节一点不增加。廉价望远镜广告里“放大六百倍”的吆喝，多半是空放大；行家先问口径，从不先问倍数。

\end{mainline}

\begin{mathbridge}
\textbf{薄透镜公式从哪里来。} 只需要图 2 和两组相似三角形。设物高 $h$，像高 $h'$。第一组：过光心的光线不偏折，物尖、光心、像尖共线，于是
\[
\frac{h'}{h}=\frac{s'}{s}.
\]
第二组：平行于主轴入射的光线过焦点。追这条光线：它在透镜平面处高度为 $h$，此后直奔焦点 $F$，再继续走到像平面时降到高度 $h'$。由焦点两侧两个相似三角形，
\[
\frac{h}{h'}=\frac{f}{s'-f}.
\]
两式相乘消去 $h'/h$：
\[
1=\frac{s'}{s}\cdot\frac{f}{s'-f}
\;\Longrightarrow\;
s(s'-f)=s'f
\;\Longrightarrow\;
ss'-sf=s'f.
\]
两边除以 $ss'f$，正是 $\dfrac{1}{f}-\dfrac{1}{s'}=\dfrac{1}{s}$，移项即成像公式。推导里偷偷用了两个假设：透镜足够薄（所有折射当作发生在同一个平面上），光线足够贴近主轴（“高”可以直接当“弧长”用）。这两个假设有多可靠，就是下一节“模型的边界”要审的案。

\textbf{屈光度为什么能相加。} 两片薄透镜紧贴时，总屈光度 $P=P_1+P_2$。理由：第一片把光锥的聚散程度改变了 $P_1$（以“米的倒数”为单位），第二片在它出来的形状上再叠加改变 $P_2$。这就是验光师把试镜片一片片往镜架上加的数学许可证——加的是屈光度，不是焦距。
\end{mathbridge}

\begin{mainline}

\step{模型的边界}

每个公式都要交代自己的辖区。光线成像理论有三道围墙。

\textbf{第一道：薄透镜与近轴。} 上面的推导默认光线都贴着主轴、以小角度行进。真实透镜口径一大，边缘入射的光线偏折角就大，近轴近似开始漏水：球面透镜的边缘光线和中心光线聚不到同一点，这叫\emph{球差}。历史上最贵的球差出在哈勃空间望远镜上：它的主镜直径 2.4 米，边缘形状与设计的偏差只有约 2 微米——不到头发丝直径的三十分之一——却足以让星光散开成模糊的光斑，只能等 1993 年航天员上天给它“配眼镜”（加装矫正光学组件）才恢复锐利。2 微米毁掉十几亿美元设备锐度的故事，值得每个觉得“近似差不多就行”的人默读三遍。

\textbf{第二道：颜色。} 玻璃对不同颜色的光折射率略有不同，所以同一片透镜对蓝光和红光的焦距不同，像的边缘会镶上彩色条纹，这叫\emph{色差}。牛顿当年被色差折磨得失去耐心，干脆绕开折射，设计了用凹面镜聚光的反射望远镜——镜子对所有颜色一视同仁。今天最大的光学望远镜无一例外是反射式，遗产就在这里。球差和色差只是像差家族里最出名的两位；斜入射的光束会拖出彗星尾巴（彗差），横竖两个方向的焦线对不齐（像散），直线条被弯成桶形或枕形（畸变）。相机镜头里动辄七八片甚至十几片镜片，大部分不是在“放大”，而是在互相抵消彼此的像差——每片透镜犯的错，由形状相反的另一片来偿。

\textbf{第三道，也是最深的：光根本不总是走直线。} 回到实验三那个反直觉的结果：遮住透镜一半，像完整无缺，只是变暗。这说明每个像点是物点发出的一整个光锥共同堆起来的，少一半光，只是每个点暗一半。这本来挺好懂，但它反过来说明：光线模型里“每一点的光各自独立走直线”的图像并不扎实。真正让围墙倒塌的是小孔。做针孔相机的人都知道，孔越小，几何上的像越锐利；可孔小到零点几毫米以下，像反而重新变模糊——光经过小孔后不肯再走直线，向四周散开。这不是透镜的缺陷，不是工艺问题，是“光走直线”这条公理本身在孔径接近光的波长（可见光约 $0.5$ 微米）时作废。第 35 章会把光当成波来重建这一切；这里先记账：针孔成像有个最佳孔径，几何模糊随孔径增大、衍射模糊随孔径减小，对十厘米深的针孔相机，最佳孔径大约在 $0.2$ 到 $0.3$ 毫米——两个趋势在真实世界里掐出的折中。

这道围墙还带来一个日常后果：\emph{分辨率有极限}。任何口径为 $D$ 的圆孔（瞳孔、镜头、主镜），点光源通过它成的像都是一个有限大小的斑，两个点光源的角距离小于大约 $\theta\sim\lambda/D$ 时就分不清谁是谁。人眼瞳孔直径约 $3$ 到 $5$ 毫米，取可见光波长 $5.5\times10^{-7}$ 米，$\theta\sim10^{-4}$ 弧度，约半角分——而正常视力表上“$1.0$”对应的恰好是分辨 $1$ 角分，视网膜黄斑处感光细胞的间距也正好排到这个密度。瞳孔的衍射极限和视网膜的“采样密度”在同一个量级上碰头，进化没有浪费任何一边。这不是巧合的赞美，而是一个提示：你眼睛的光学设计已经顶到了物理允许的天花板，想看得更细，只能换更大的口径（望远镜）或更短的波长。

\end{mainline}

\begin{challenge}
\textbf{瑞利判据与那个 1.22。} 更精确的分辨极限是 $\theta_{\min}\approx1.22\lambda/D$。系数 $1.22$ 来自圆孔衍射图样的第一个暗环位置——圆孔的衍射图样由贝塞尔函数描述，第一零点在 $x\approx3.83$ 处，而 $3.83/\pi\approx1.22$。推导本身属于第 35 章的波动光学，这里只强调逻辑地位：$1.22$ 不是拟合出来的经验数，是“波通过圆孔”这个几何问题的确定答案。对显微镜，相应的结果是横向分辨极限 $d\sim\lambda/(2\,\mathrm{NA})$，其中数值孔径 $\mathrm{NA}=n\sin\theta$ 描述物镜收集光锥的胖瘦。油浸物镜用油把 $n$ 从 $1$ 提到 $1.5$ 左右，就是在分子上做文章。可见光显微镜的分辨极限停在约 $0.2$ 微米——比这个更小的东西，不是“放大倍数不够”，而是可见光这个探针本身太粗。电子显微镜之所以能看病毒和晶格，是因为加速电子的德布罗意波长可以短到皮米量级；光学显微镜的天花板，正是后面量子篇“物质也是波”的第一块敲门砖。

\textbf{反射式望远镜的口径账。} 把 $\theta\sim\lambda/D$ 反过来用：想分辨角距离 $\theta$ 的两个天体，$D$ 至少要 $\lambda/\theta$。这也是为什么射电望远镜要造到几百米——不是可见光望远镜的放大野心，而是波长厘米级的波想达到同样分辨率时必须付的口径税。
\end{challenge}

\begin{mainline}

\step{回头解释开场现象}

现在批卷子，从最难缠的“左右颠倒”开始。

盯住图 1：镜子做过的事只有一件——把每条光线按反射定律折回。设镜面竖直，物点 $P$ 与像点 $P'$ 的连线\emph{垂直于镜面}，到镜面等距。也就是说，镜子反转的方向既不是左右也不是上下，而是\emph{前后}：沿着视线看进去的那个方向被翻了过来。你的鼻尖离镜面 30 厘米，镜中鼻尖就在镜后 30 厘米；你的后背离镜面 30.4 厘米，镜中后背就在镜后 30.4 厘米——前后次序整个倒转。那“左右颠倒”的错觉从哪来？来自你做比较的方式：你想象自己绕到镜子后面、\emph{原地转身}面对现在的自己，然后发现举左手对应的是他的右手。问题出在这个想象动作上——你是绕着竖直轴转的身。如果你改成绕水平轴翻个跟头（头朝下脚朝上）去跟镜中人比，你会发现被“颠倒”的是上下，左右反而一致。镜中人只是在“前后”这个方向上是你的反面；“左右颠倒”还是“上下颠倒”，取决于你假设他用了哪种方式转过身来。镜子对水平和竖直一视同仁，是你的转身方式偏心了竖直轴。

再看半身镜。从你的眼睛向镜中虚像的头顶连一条直线——那就是实际反射光线从镜面出发的方向；再向虚像的脚底连一条。这两条线在镜面上截出的那段长度，由相似三角形可知恰好是身高的一半，而且这个比例里根本不出现你离镜子的距离：你退后，虚像也等比例退后，两条线在镜面处的间距不变。便利贴实验不动的结果，就是这么来的。镜子只要够身高的一半、挂在正确的高度，全身永远装得下。

汽车后视镜那句警告也已结案：凸面镜 $f<0$，永远成正立缩小虚像，大脑按“像的大小”估计距离，于是把后车判远了。工程师用缩小换来了大视野，并在镜面上印字，替你的直觉交学费。

汤匙实验也顺便销案。汤匙凹面是一面小凹面镜，焦距只有几厘米。你把脸凑到焦距以内，满足 $s<f$，看到正立放大的虚像；拿远到焦点附近，反射光趋向平行，像糊成一片；再远，$s>f$，光线真的会聚了，成倒立实像——只是这个实像悬浮在你和汤匙之间的空气里，你的眼睛位于它后面，把它当远处的东西接收下来。凸面 $f<0$，从公式看 $s'$ 永远为负、$m$ 永远在 $0$ 和 $1$ 之间：正立缩小，童叟无欺。一把餐具走完了整个成像公式。

还剩一个小赠品。镜中的你是“过去”的你：光从你到镜子再回来，若你离镜 1 米，光走了 2 米，约 $7$ 纳秒。你和镜中人对视，看到的永远是对方 $7$ 纳秒前的样子。这个延迟小到荒谬，但它属于一个大的事实：任何“看见”都不是即时的，都是光线花了时间、走了路径之后的结算单。

\step{脑洞实验与挑战题}

\begin{thought}[给月球上的国旗拍张照]

常有人说“望远镜厉害到能看见宇航员插在月球上的国旗”。我们用本章的极限公式给它称称重。国旗展开约 1 米，月地距离约 $3.8\times10^{8}$ 米，它对地球的张角约 $2.6\times10^{-9}$ 弧度。要分辨这个张角，口径至少 $D\sim\lambda/\theta\approx5.5\times10^{-7}/2.6\times10^{-9}\approx200$ 米——注意这还是没乘 $1.22$ 的乐观账。目前最大的光学望远镜口径 10 米量级，差了二十多倍；哈勃更不行。所以没人拍到过那面国旗，未来也不会用单面镜子拍到。但反过来想：这个限制来自波长和口径之比，不是来自“技术还不够好”。人类的确“看”到了登月遗迹——是用月球勘测轨道飞行器在约 50 公里高的环月轨道上拍的，把物距从 $3.8\times10^{8}$ 米缩到 $5\times10^{4}$ 米，需要的口径立刻掉到几厘米。物理判了死刑的案子，工程师换了个法庭重审：打不过衍射极限，就搬得离它近一点。

这个脑洞还有一层值得翻出来看：衍射极限是“光以波动方式传播”的证据，却在几何光学的章节里提前收了税。这说明“光线”模型和“波动”模型不是两个可以随便挑的世界观——前者是后者在“孔径远大于波长”时的近似。本章所有公式都活在这个近似里，而它失效的方式（小孔散开、像点成斑）恰恰泄露了下一章的真相。
\end{thought}

\begin{puzzles}
\item 潜望镜用两块互相平行的平面镜，把光折两次送进眼睛。隔着潜望镜看一张写着字的纸，字是正的还是反的？如果三块镜子折三次呢？（提示：每照一次平面镜，前后方向反转一次；数清楚反转次数，再想想偶数次反转后“手性”回到什么状态。）
\item 把放大镜完全浸在水里看池底的石子，放大效果比在空气里强还是弱？（提示：透镜的会聚本事来自“玻璃和周围介质的折射率之差”，水的折射率 1.33 比空气 1.00 更接近玻璃，焦距会怎么变？再想想鱼的晶状体为什么比人的凸得多。）
\item 实验三里，遮住凸透镜的一半，像变暗但保持完整。现在换一个问题：把透镜的\emph{上半部分}换成另一片焦距稍长的透镜（下半不动），纸上会发生什么？（提示：回到“每个像点由整个光锥堆成”这句话；现在的光锥被劈成了两股各自会聚的束。）
\end{puzzles}

最后，用全书反复追问的三个问题收束本章。系统此刻是什么状态？一组器件（镜、透镜、眼球）加一组光锥，每个物点的光锥此刻被改造成了会聚或发散的某种形状。什么规律决定状态如何变化？反射定律与折射定律——在第 33 章里它们又被费马原理统一成“光取传播时间取驻值的路”——加上一个把它们压缩到极限的成像方程。我们怎样通过测量知道它？把白纸伸到像的位置看光斑是否真实会聚，在镜面上贴便利贴量距离，遮住透镜的一半看像是否残缺。下一章我们将看到，连“光走直线”这条最老的公理，也只是在某种尺度下足够好的近似——物理学的每面镜子背后，都还有一面镜子。

\end{mainline}
